\uuid{9TbH}
\titre{Polynôme irréductible}
\auteur{bodin ; hanani ; mzari}
\organisation{exo7}
\licence{CC BY-NC-SA 4.0}
\datecreate{2020-01-20}
\niveau{L1}
\module{Algèbre}
\chapitre{Polynôme, fraction rationnelle}
\sousChapitre{Racine, décomposition en facteurs irréductibles}
\difficulte{3}
\theme{Racines, factorisation, polynôme irréductible}
\exo7qcmid{lille-1.polynomes.023}

\contenu{
\begin{question}
\qclassification{105.03}
\qtags{motcle=polynôme irréductible}
Soit $P \in \Qq[X]$ un polynôme de degré $n$.
\begin{answers}
    \good{$P$ peut avoir des racines dans $\Rr$, mais pas dans $\Qq$.}

    \good{Si $z\in \Cc\setminus\Rr$ est une racine de $P$, alors $\bar z$ aussi.}

    \bad{Les facteurs irréductibles de $P$ sur $\Qq$ sont de degré $1$ ou $2$.}

    \bad{Les racines réelles de $P$ sont de la forme $\alpha + \beta\sqrt{\gamma}$, $\alpha,\beta,\gamma \in \Qq$.}
\end{answers}
\begin{explanations}
Sur $\Qq$ les facteurs irréductibles peuvent être de n'importe quel degré.
\end{explanations}
\end{question}
}
